% Učitelský klíč k pracovnímu listu 64 (nikoli žákovská verze).
\documentclass[11pt,a4paper]{article}
\input{../../../sablony/pripravy/priprava.sty}
\nastavPripravu{1. ročník}{Funkce — řešení PL}{64}{Lineární funkce}
\begin{document}
\blok{Úloha 1 — grafy z předpisů}
Pro konstrukci stačí dva body:
\[
\begin{array}{c|cc|cc}
 & a & b & [x_1;y_1] & [x_2;y_2]\\\hline
f(x)=x-2 & 1 & -2 & [0;-2] & [2;0]\\
g(x)=-2x+1 & -2 & 1 & [0;1] & [1;-1]\\
h(x)=\frac12x+2 & \frac12 & 2 & [0;2] & [2;3]
\end{array}
\]
Všechny tři grafy jsou při definičním oboru $\mathbb R$ přímky.
\blok{Úloha 2 — parametry a monotonie}
$f$ roste, $g$ klesá, $h$ roste. Průsečíky s osou $y$ jsou
$[0;-2]$, $[0;1]$ a $[0;2]$.
\blok{Úloha 3 — předpisy ze dvou bodů}
\begin{enumerate}[label=\alph*)]
\item $A[0;2]$, $B[2;6]$: $a=\frac{6-2}{2-0}=2$, $b=2$,
      tedy $f(x)=2x+2$. Zkouška $f(2)=6$.
\item $C[-2;0]$, $D[2;-2]$: $a=\frac{-2-0}{2-(-2)}=-\frac12$.
      Rovnice $0=(-\frac12)(-2)+b$ dává $b=-1$;
      tedy $g(x)=-\frac12x-1$. Zkouška $g(2)=-2$.
\end{enumerate}
\blok{Úloha 4 — předpisy z grafu}
\begin{itemize}
\item Plná $p$: body $[0;-1]$ a $[1;1]$, tedy $a=2$, $b=-1$ a $p(x)=2x-1$.
\item Čárkovaná $q$: body $[0;3]$ a $[1;2]$, tedy $a=-1$, $b=3$ a $q(x)=-x+3$.
\item Tečkovaná $r$: body $[0;1]$ a $[2;2]$, tedy $a=\frac12$, $b=1$ a $r(x)=\frac12x+1$.
\end{itemize}
\blok{Úloha 5 — posun}
Grafy $y=x$, $y=x+2$ a $y=x-1$ jsou navzájem rovnoběžné, všechny mají směrnici $1$.
Absolutní člen $b$ způsobí svislý posun o $b$ vzhledem k přímce $y=x$.
\blok{Úloha 6 — konečná množina parametru}
Pro $a\in\{-2;0;1\}$ vznikají právě tři předpisy:
\[
 y=-2x+1,\qquad y=1,\qquad y=x+1.
\]
Všechny přímky procházejí bodem $[0;1]$; mají odlišné směrnice.
\blok{Úloha 7 — parametr z intervalu}
Krajní přímky jsou $y=x-1$ a $y=x+2$. Pro libovolné reálné
$b\in\langle-1;2\rangle$ je to jiná přímka se směrnicí $1$.
Je jich nekonečně mnoho a tvoří rodinu funkcí, nikoli jedinou funkci.
\blok{Úloha 8 — omezený definiční obor}
$D(f)=\langle-2;1)$, neboť tak zní zadání. Koncové body grafu jsou
$[-2;4]$ (plný) a $[1;1]$ (prázdný). Funkce klesá, takže
$H(f)=(1;4\rangle$.
\blok{Úloha 9 — bonus}
$a=(7-3)/(3-1)=2$, z $3=2+b$ plyne $b=1$, proto $y=2x+1$.
Graf je rovnoběžný s $y=2x-1$ a leží o dvě jednotky výše.
\poznamka{U uznání grafů trvej na zřejmém měřítku, správně zakreslených bodech
 a jednoznačném přiřazení předpisu. U úlohy 7 lze místo jednotlivých přímek názorně vyznačit
 pás, ale vysvětli, že pás není grafem jedné funkce.}
\end{document}
